\section{Le bilinguisme des langues officielles canadiennes}\label{le-bilinguisme-des-langues-officielles-canadiennes}

\stanceinfo{\textbf{CCLI 2024} [2024-01-05]}{RdP2024 [2024-09-22]}

\officialstance{La Fédération canadienne étudiante de génie soutient les deux langues officielles du Canada, le français et l'anglais, conformément à la Charte canadienne des droits et libertés. En tant qu'organisme à but non lucratif, la Fédération canadienne étudiante de génie accorde la plus haute priorité au bilinguisme. Il s'agit de réduire l'écart systémique entre les membres francophones et anglophones de l'association, afin que les opportunités ne soient pas limitées par des barrières linguistiques.}

\stanceheading{La position des étudiants}
\begin{itemize}
 \item Il existe actuellement une barrière linguistique entre les membres anglophones et francophones.
 \item Il n'y a pas de procédures normalisées pour les conférences et les activités entièrement bilingues de la FCEG.
 \item Il y a un manque d'universités bilingues à travers le Canada, ainsi qu'un manque de ressources pour apprendre à parler couramment l'une ou l'autre des langues officielles dans le cadre du diplôme d'ingénierie.
 \item La FCEG estime que tous les membres devraient se sentir respectés, confiants et entendus lorsqu'ils s'expriment ou dirigent les activités de la FCEG dans l'une ou l'autre langue.
 \item L'augmentation du bilinguisme conduit à un espace plus inclusif, avec la possibilité de communiquer plus efficacement entre les membres.
\end{itemize}

\stanceheading{Les recommandations aux partenaires, aux parties intéressées et aux autres entités}
\begin{itemize}
 \item La FCEG demande au BCAPG de revoir la politique de choix de la langue afin d'offrir l'option d'apprendre la deuxième langue officielle du Canada, si l'on ne connaît pas les deux langues.
 \item La FCEG demande au BCAPG de collaborer avec les DDIC afin de trouver une place pour le bilinguisme dans toutes les écoles, quelle que soit la langue majoritairement parlée dans un établissement.
 \item La FCEG demande aux DDIC de soutenir l'enseignement bilingue et d'en faire l'une de leurs priorités stratégiques à l'avenir afin de combler le fossé linguistique et d'offrir davantage de possibilités aux étudiants en ingénierie au niveau national.
 \item La FCEG recherchera des partenariats pour des services d'éducation des adultes qui offrent des possibilités d'éducation en français et en anglais (tutorat, cours, etc.) aux membres de la FCEG à un tarif réduit. L'objectif est d'encourager l'apprentissage extrascolaire facultatif afin de surmonter la barrière linguistique entre les étudiants francophones et anglophones.
 \item La FCEG demande à ses partenaires de fournir des présentations bilingues afin que les anglophones et les francophones puissent mieux comprendre et suivre ce qui est présenté.
 \item La FCEG appelle les responsables d'activités et les officiers de la FCEG à s'assurer que toutes les activités de la FCEG restent des espaces sûrs pour les minorités francophones grâce à l'utilisation de traductions précises et d'une interprétation en direct.
 \item La FCEG appelle les membres à rechercher des opportunités pour promouvoir le bilinguisme parmi les étudiants et à créer des initiatives qui soutiennent le bilinguisme dans leurs institutions.
\end{itemize}

\stanceheading{Ce que fait la FCEG}
\begin{itemize}
 \item Utiliser l'interprétation en direct lors des assemblées générales et d'autres réunions politiques officielles au sein de la FCEG.
 \item Enquêter sur la politique en matière de choix de langues au sein du BCAPG afin de déterminer quelles universités limitent le nombre de cours de l'autre langue officielle qui peuvent être suivis.
 \item Utiliser les outils actuels d'intelligence artificielle pour faciliter la traduction en ligne et en personne.
 \item La FCEG dispose d'un commissaire au bilinguisme, supervisé par le vice-président des communications, ainsi que d'un groupe de travail sur le bilinguisme chargé d'aider à la traduction des documents, des présentations et des politiques.
 \item Lors de l'élection des exécutifs nationaux de la FCEG, chaque question doit être posée dans les deux langues officielles en alternant la langue posée en premier. Tous les candidats seront interrogés sur leur niveau de compréhension du français et de l'anglais ainsi que sur leur plan de communication avec les délégués dont la langue maternelle ne correspond pas à la leur.
\end{itemize}

\stanceheading{Ce que la FCEG envisage de faire}
\begin{itemize}
 \item Créer une procédure normalisée pour la traduction et l'interprétation et planifier des procédures spécifiques pour les activités de la FCEG afin de s'assurer que toutes les activités et conférences sont entièrement bilingues et constituent une priorité dès le départ.
 \item Soutenir les agents de la FCEG qui ne maîtrisent pas les deux langues officielles en leur fournissant des ressources pour les aider dans la langue dans laquelle ils ne sont pas à l'aise.
 \item Veiller à créer des espaces sûrs où les deux langues peuvent être parlées librement, avec l'aide d'un interprète si nécessaire.
 \item Donner la priorité aux présentateurs bilingues dans les événements et les activités des conférences de la FCEG, en prévoyant des alternatives pour les présentateurs unilingues afin qu'ils puissent être compris par les participants à la conférence.
 \item Mettre en œuvre un formulaire de pratique du bilinguisme qui permette de recueillir les commentaires des membres.
 \item Étudier la possibilité d'un partenariat avec des services de traduction et des organisations bilingues spécialisées dans le domaine de l'ingénierie.
 \item Mettre en place un système de mentorat formel et informel au sein de la FCEG afin d'encourager la pratique des langues dans un environnement peu stressant.
 \item Créer une boîte à outils du bilinguisme de la FCEG à l'intention des membres, qui comprendra notamment les éléments suivants
 \begin{itemize}
  \item une formation au bilinguisme;
  \item un modèle pour les présentations bilingues;
  \item les meilleures pratiques en matière d'intercommunication;
  \item une liste de mots et d'expressions;
  \item des mécanismes de rétroaction.
 \end{itemize}
\end{itemize}

\newpage
